% Exercice: Exercice 2
% Sujet: Logique
% Généré automatiquement

\begin{exercice}[Logique]
\label{ex:4}

Exercice 2 $(4.5 points)$ _\\
Après un contrôle les notes de mathématiques de 6^\{0\} élèves dune classe de 1èrD\\
ont été regroupées dans le tableau suivant\\
Notes\\
[0; [\\
[1 ;8[\\
[8,1^\{2\}[\\
[1^\{2\} ;1^\{6\}[\\
[1^\{6\} ;2^\{0\}[\\
Nombre délèves\\
Fréquences\\
0,3\\
Effectifs cumulés croissants\\
[2pts]\\
Recopier et compléter ce tableau .\\
Calculer le pourcentage des élèves ayant une note supérieure ou égale\\
1^\{2\}/2^\{0\}\\
[0.Spt]\\
Lépreuve de mathématiques du contrôle est constituée de\\
exercices et dun problème.\\
Lenseignant dispose dans sa banque d'épreuve de 1^\{8\} exercices (dont\\
9 en\\
en statistiques\\
trigonométrie) et 1^\{0\} problèmes différents.\\
équations\\
Combien d'épreuves différentes peut-\\
[Ipt]\\
composer\\
Donner le nombre dépreuves contenant exactement un exercice de statistique et un exercice\\
de trigonométrie\\
[lpt] :\\

\end{exercice}

% --- Fin Exercice 2 ---
